\documentclass[aspectratio=169,11pt]{beamer}
\usetheme{metropolis}
\usepackage[T1]{fontenc}
\usepackage{lmodern,microtype}
\usepackage{amsmath,amssymb}
\usepackage{booktabs}
\usepackage{xcolor}
\usepackage[numbers,sort&compress]{natbib}
\usepackage{url}
\usepackage{hyperref}
\hypersetup{colorlinks=true,urlcolor=blue,citecolor=blue,linkcolor=blue}

\title{How Models Represent, Store, and Retrieve Knowledge}
\subtitle{Module Q13 --- DLS 2026}
\author{Deep Learning Systems}
\date{\today}

\begin{document}

\maketitle

\begin{frame}{The student's question, sharpened}
  \emph{``How do models represent and store knowledge, and use it to answer?''}
  \begin{itemize}
    \item We separate three concerns:
    \begin{enumerate}
      \item \textbf{Representation}: how facts become geometry in activation/weight space.
      \item \textbf{Storage}: which parameters hold what, and how retrieval happens.
      \item \textbf{Use}: the forward-pass mechanics that turn a prompt into an answer.
    \end{enumerate}
    \item Key tension: knowledge is \emph{distributed} (spread over many weights) yet
      often \emph{localizable} (editable at specific sites)~\cite{meng2022rome}.
    \item We treat the transformer as a read--write system over a shared residual stream~\cite{elhage2021mathematical}.
  \end{itemize}
\end{frame}

\begin{frame}{Representation: the distributional and linear view}
  \begin{itemize}
    \item \textbf{Distributional hypothesis}: meaning is inferred from co-occurrence;
      embeddings place similar tokens near one another~\cite{mikolov2013efficient}.
    \item \textbf{Linear structure}: relations appear as approximately constant offset
      vectors (analogies), evidence that attributes are directions, not slots~\cite{mikolov2013linguistic}.
    \item \textbf{Superposition}: a $d$-dimensional space stores more features than
      dimensions by packing them in near-orthogonal directions~\cite{elhage2022superposition}.
    \item Consequence: a single neuron is usually \emph{polysemantic}; the natural unit
      of knowledge is a \emph{direction} in the residual stream, recoverable via sparse
      dictionaries~\cite{bricken2023monosemanticity}.
  \end{itemize}
\end{frame}

\begin{frame}{Storage: feed-forward layers as key--value memories}
  \begin{itemize}
    \item Each MLP block computes $\text{FFN}(x)=\sigma(x W_1)\,W_2$.
    \item Interpret rows of $W_1$ as \textbf{keys} (patterns they fire on) and rows of
      $W_2$ as \textbf{values} (what they write back)~\cite{geva2021kvmemories}.
    \item Empirically, FFN layers act as a differentiable associative memory: a key
      matches a context, its value injects a prediction into the residual stream.
    \item Factual associations are concentrated in mid-layer FFNs at the final token of
      the subject, enabling causal localization~\cite{meng2022rome}.
    \item Hence ``storage'' $\approx$ a learned hash of context $\to$ output contribution.
  \end{itemize}
\end{frame}

\begin{frame}{Use: moving and combining information with attention}
  \begin{itemize}
    \item The residual stream is a communication bus; attention \emph{moves} information
      between positions, FFNs \emph{process} it~\cite{elhage2021mathematical}.
    \item \textbf{Induction heads} implement in-context copying/pattern completion and
      correlate with the emergence of in-context learning~\cite{olsson2022incontext}.
    \item A simple factual-recall circuit: attention gathers the subject representation
      $\to$ mid-layer FFN keys match $\to$ value writes the object $\to$ unembedding reads
      the answer logit.
    \item \textbf{In-context knowledge} (from the prompt) and \textbf{parametric knowledge}
      (from weights) compose in the same stream; conflicts are resolved by relative
      magnitude of their contributions.
  \end{itemize}
\end{frame}

\begin{frame}{Evidence it is localizable: probing and editing}
  \begin{itemize}
    \item \textbf{Causal tracing}: corrupt the subject, restore single hidden states, and
      measure recovery to pinpoint where a fact lives~\cite{meng2022rome}.
    \item \textbf{Rank-one editing (ROME)}: write a new value into the located FFN to change
      a specific fact while leaving others intact~\cite{meng2022rome}.
    \item \textbf{Mass editing (MEMIT)}: scale editing to thousands of facts across layers~\cite{meng2022memit}.
    \item Caveat: edits can be brittle and leak---specificity and generalization trade off,
      so ``localized'' does not mean ``cleanly isolated''.
  \end{itemize}
\end{frame}

\begin{frame}{When parametric memory is not enough: retrieval augmentation}
  \begin{itemize}
    \item Parametric storage is finite, stale, and hard to attribute.
    \item \textbf{RAG}: condition generation on documents fetched by a learned retriever,
      decoupling knowledge from parameters~\cite{lewis2020rag}.
    \item \textbf{Retrieval during pretraining} (RETRO) lets a smaller model match larger
      ones by reading a trillion-token datastore~\cite{borgeaud2022retro}.
    \item Benefits: updatable knowledge, source attribution, reduced hallucination on
      long-tail facts; cost: retriever quality and context-length pressure.
  \end{itemize}
\end{frame}

\begin{frame}{Failure modes: why stored knowledge misleads}
  \begin{itemize}
    \item \textbf{Hallucination}: fluent output ungrounded in stored or retrieved facts;
      a surveyed, systemic property of generation~\cite{ji2023hallucination}.
    \item \textbf{Superposition interference}: packed features can fire spuriously,
      producing confident errors~\cite{elhage2022superposition}.
    \item \textbf{Stale parametric facts}: weights freeze the training snapshot; editing
      and retrieval are the two remedies above.
    \item \textbf{Attribution gap}: distributed storage makes it hard to say \emph{why}
      a model answered---motivating interpretability tooling.
  \end{itemize}
\end{frame}

\begin{frame}{Synthesis and takeaways}
  \begin{itemize}
    \item \textbf{Represent}: knowledge is directions in a superposed activation space,
      not discrete symbols~\cite{elhage2022superposition,bricken2023monosemanticity}.
    \item \textbf{Store}: FFNs act as key--value memories; facts are distributed yet
      partly localizable and editable~\cite{geva2021kvmemories,meng2022rome}.
    \item \textbf{Use}: attention routes information, FFNs recall it, the unembedding reads
      it out as logits~\cite{elhage2021mathematical,olsson2022incontext}.
    \item \textbf{Engineer}: when parametric memory fails, retrieve
      (RAG/RETRO) or edit (ROME/MEMIT)~\cite{lewis2020rag,borgeaud2022retro,meng2022memit}.
  \end{itemize}
\end{frame}

\begin{frame}[allowframebreaks]{References}
  \footnotesize
  \begin{thebibliography}{99}
    \bibitem{mikolov2013efficient} T.~Mikolov, K.~Chen, G.~Corrado, and J.~Dean.
      Efficient Estimation of Word Representations in Vector Space. \emph{arXiv:1301.3781}, 2013.
    \bibitem{mikolov2013linguistic} T.~Mikolov, W.~Yih, and G.~Zweig.
      Linguistic Regularities in Continuous Space Word Representations. \emph{NAACL-HLT}, 2013.
    \bibitem{elhage2021mathematical} N.~Elhage et al.
      A Mathematical Framework for Transformer Circuits. \emph{Anthropic / Transformer Circuits Thread}, 2021.
    \bibitem{elhage2022superposition} N.~Elhage et al.
      Toy Models of Superposition. \emph{Transformer Circuits Thread}, 2022.
    \bibitem{bricken2023monosemanticity} T.~Bricken et al.
      Towards Monosemanticity: Decomposing Language Models With Dictionary Learning. \emph{Transformer Circuits Thread}, 2023.
    \bibitem{geva2021kvmemories} M.~Geva, R.~Schuster, J.~Berant, and O.~Levy.
      Transformer Feed-Forward Layers Are Key-Value Memories. \emph{EMNLP}, 2021.
    \bibitem{meng2022rome} K.~Meng, D.~Bau, A.~Andonian, and Y.~Belinkov.
      Locating and Editing Factual Associations in GPT. \emph{NeurIPS}, 2022.
    \bibitem{meng2022memit} K.~Meng, A.~Sharma, A.~Andonian, Y.~Belinkov, and D.~Bau.
      Mass-Editing Memory in a Transformer. \emph{ICLR}, 2023.
    \bibitem{olsson2022incontext} C.~Olsson et al.
      In-context Learning and Induction Heads. \emph{Transformer Circuits Thread}, 2022.
    \bibitem{lewis2020rag} P.~Lewis et al.
      Retrieval-Augmented Generation for Knowledge-Intensive NLP Tasks. \emph{NeurIPS}, 2020.
    \bibitem{borgeaud2022retro} S.~Borgeaud et al.
      Improving Language Models by Retrieving from Trillions of Tokens. \emph{ICML}, 2022.
    \bibitem{ji2023hallucination} Z.~Ji et al.
      Survey of Hallucination in Natural Language Generation. \emph{ACM Computing Surveys}, 2023.
  \end{thebibliography}
\end{frame}

\end{document}